% GENERATED FILE. Edit canonical lessons, then run npm run generate:books.
% canonical-source-hash: fnv1a64:60947ede14d9b387
% canonical-lessons: RU-C394-dolg, RU-C394-kredit, RU-C394-odalzhivat, RU-C394-vozvrashchat, RU-C394-procent

\chapter{Debts and loans}
\label{ch:ru-debts-and-loans}
\hlchaptermodality{394}

\begin{chapteropening}
\textbf{By the end of this chapter:} \emph{I can talk about a debt and a loan, lend and borrow money, pay it back and say what interest a bank pays.}

Complete the last lesson of chapter 394: I can talk about a debt and a loan, lend and borrow money, pay it back and say what interest a bank pays.
\end{chapteropening}

\section[\texorpdfstring{\ru{долг} (dolg)}{долг (dolg)}]{\ru{долг} (dolg) — a debt; a duty}

\label{lesson:RU-C394-dolg}



\begin{quote}
Before the new one: say the Russian for a power of attorney, then the Russian for a will.
\end{quote}

\subsection*{You'll want to know: \ru{долг}}
\textbf{\ru{долг}} (\emph{dolg}) — \textquotedblleft{}a debt; a duty\textquotedblright{}.

\begin{quote}
\ru{У} \ru{него} \ru{большие} \ru{долги}.
\end{quote}

\begin{quote}
\emph{U negó bolshíye dolgí.}
\end{quote}

\begin{quote}
He has big debts.
\end{quote}

Masculine; plural \emph{\ru{долги}}. \emph{\ru{Брать} \ru{в} \ru{долг}} — to borrow. \emph{\ru{Долг}} is also a duty: \emph{\ru{мой} \ru{долг} — \ru{помочь}}.

\begin{cousinweb}
\textbf{\ru{долг}} is inherited from Proto-Slavic \emph{*dŭlgŭ} (EtymDB) \textquotedblleft{}a debt, an obligation\textquotedblright{}.
\end{cousinweb}

\begin{grammarlens}[title={What you've built}]
One more word for telling, explaining and describing at B1.
\end{grammarlens}

\subsection*{Your turn}
\begin{itemize}
\raggedright
  \item \emph{Say it:} \emph{dolg}
  \item \emph{Say it:} \emph{dolg}, once more
  \item \emph{Say it:} a sentence of your own with \emph{dolg}
  \item \emph{Recall:} say the Russian for a power of attorney, then the Russian for a will, then say \emph{dolg} again
\end{itemize}

\begin{tcolorbox}[breakable,colback=teal!4,colframe=teal!35!black,title={Before you move on}]
What does \ru{долг} mean? (\textquotedblleft{}A debt; a duty\textquotedblright{}.) Say it once more.
\end{tcolorbox}
\section[\texorpdfstring{\ru{кредит} (kredít)}{кредит (kredít)}]{\ru{кредит} (kredít) — a loan, credit}

\label{lesson:RU-C394-kredit}



\begin{quote}
Before the new one: say the Russian for a will, then the Russian for a debt.
\end{quote}

\subsection*{You'll want to know: \ru{кредит}}
\textbf{\ru{кредит}} (\emph{kredít}) — \textquotedblleft{}a loan, credit\textquotedblright{}.

\begin{quote}
\ru{Они} \ru{взяли} \ru{кредит} \ru{на} \ru{машину}.
\end{quote}

\begin{quote}
\emph{Oní vzyáli kredít na mashínu.}
\end{quote}

\begin{quote}
They took out a loan for a car.
\end{quote}

Masculine. \emph{\ru{Брать} \ru{кредит}} — to take out a loan. \emph{\ru{В} \ru{кредит}} — on credit.

\begin{cousinweb}
\textbf{\ru{кредит}} is usually traced, through German and French, to Latin \emph{crēditum} \textquotedblleft{}a loan\textquotedblright{}, from \emph{crēdere} \textquotedblleft{}to trust, to lend\textquotedblright{} (Lewis \& Short): a loan rests on trust.
\end{cousinweb}

\begin{grammarlens}[title={What you've built}]
One more word for telling, explaining and describing at B1.
\end{grammarlens}

\subsection*{Your turn}
\begin{itemize}
\raggedright
  \item \emph{Say it:} \emph{kredít}
  \item \emph{Say it:} \emph{kredít}, once more
  \item \emph{Say it:} a sentence of your own with \emph{kredít}
  \item \emph{Recall:} say the Russian for a will, then the Russian for a debt, then say \emph{kredít} again
\end{itemize}

\begin{tcolorbox}[breakable,colback=teal!4,colframe=teal!35!black,title={Before you move on}]
What does \ru{кредит} mean? (\textquotedblleft{}A loan, credit\textquotedblright{}.) Say it once more.
\end{tcolorbox}
\section[\texorpdfstring{\ru{одалживать} (odálzhivat)}{одалживать (odálzhivat)}]{\ru{одалживать} (odálzhivat) — to lend; to borrow}

\label{lesson:RU-C394-odalzhivat}



\begin{quote}
Before the new one: say the Russian for a debt, then the Russian for a loan.
\end{quote}

\subsection*{You'll want to know: \ru{одалживать}}
\textbf{\ru{одалживать}} (\emph{odálzhivat}) — \textquotedblleft{}to lend; to borrow\textquotedblright{}.

\begin{quote}
\ru{Я} \ru{часто} \ru{одалживаю} \ru{брату} \ru{деньги}.
\end{quote}

\begin{quote}
\emph{Ya chásto odálzhivayu brátu déngi.}
\end{quote}

\begin{quote}
I often lend my brother money.
\end{quote}

Imperfective; perfective \emph{\ru{одолжить}}. \emph{\ru{Одалживать} \ru{кому}} — to lend to; \emph{\ru{одалживать} \ru{у} \ru{кого}} — to borrow from.

\begin{cousinweb}
\textbf{\ru{одалживать}} is the imperfective of \emph{\ru{одолжить}}, built on \emph{\ru{долг}} \textquotedblleft{}a debt\textquotedblright{}, Proto-Slavic \emph{*dŭlgŭ} (EtymDB): to put someone in your debt. The \emph{\ru{о}} turns to \emph{\ru{а}} in the imperfective, as in \emph{\ru{спрашивать}}.
\end{cousinweb}

\begin{grammarlens}[title={What you've built}]
One more word for telling, explaining and describing at B1.
\end{grammarlens}

\subsection*{Your turn}
\begin{itemize}
\raggedright
  \item \emph{Say it:} \emph{odálzhivat}
  \item \emph{Say it:} \emph{odálzhivat}, once more
  \item \emph{Say it:} a sentence of your own with \emph{odálzhivat}
  \item \emph{Recall:} say the Russian for a debt, then the Russian for a loan, then say \emph{odálzhivat} again
\end{itemize}

\begin{tcolorbox}[breakable,colback=teal!4,colframe=teal!35!black,title={Before you move on}]
What does \ru{одалживать} mean? (\textquotedblleft{}To lend; to borrow\textquotedblright{}.) Say it once more.
\end{tcolorbox}
\section[\texorpdfstring{\ru{возвращать} (vozvrashchát)}{возвращать (vozvrashchát)}]{\ru{возвращать} (vozvrashchát) — to give back, to return (something)}

\label{lesson:RU-C394-vozvrashchat}



\begin{quote}
Before the new one: say the Russian for a loan, then the Russian for to lend.
\end{quote}

\subsection*{You'll want to know: \ru{возвращать}}
\textbf{\ru{возвращать}} (\emph{vozvrashchát}) — \textquotedblleft{}to give back, to return (something)\textquotedblright{}.

\begin{quote}
\ru{Он} \ru{всегда} \ru{возвращает} \ru{книги} \ru{вовремя}.
\end{quote}

\begin{quote}
\emph{On vsegdá vozvrashcháyet knígi vovrémya.}
\end{quote}

\begin{quote}
He always returns books on time.
\end{quote}

Imperfective; perfective \emph{\ru{вернуть}} or \emph{\ru{возвратить}}. \emph{\ru{Возвращать} \ru{долг}} — to pay back a debt. \emph{\ru{Возвращаться}} is to come back.

\begin{cousinweb}
\textbf{\ru{возвращать}} is built from \emph{\ru{воз}-} \textquotedblleft{}back\textquotedblright{} and \emph{\ru{вращать}} \textquotedblleft{}to turn\textquotedblright{}, the Church Slavonic twin of native \emph{\ru{ворочать}}: to turn a thing back to its owner.
\end{cousinweb}

\begin{grammarlens}[title={What you've built}]
One more word for telling, explaining and describing at B1.
\end{grammarlens}

\subsection*{Your turn}
\begin{itemize}
\raggedright
  \item \emph{Say it:} \emph{vozvrashchát}
  \item \emph{Say it:} \emph{vozvrashchát}, once more
  \item \emph{Say it:} a sentence of your own with \emph{vozvrashchát}
  \item \emph{Recall:} say the Russian for a loan, then the Russian for to lend, then say \emph{vozvrashchát} again
\end{itemize}

\begin{tcolorbox}[breakable,colback=teal!4,colframe=teal!35!black,title={Before you move on}]
What does \ru{возвращать} mean? (\textquotedblleft{}To give back, to return (something)\textquotedblright{}.) Say it once more.
\end{tcolorbox}
\section[\texorpdfstring{\ru{процент} (protsént)}{процент (protsént)}]{\ru{процент} (protsént) — a per cent; interest (on money)}

\label{lesson:RU-C394-procent}



\begin{quote}
Before the new one: say the Russian for to lend, then the Russian for to give back.
\end{quote}

\subsection*{You'll want to know: \ru{процент}}
\textbf{\ru{процент}} (\emph{protsént}) — \textquotedblleft{}a per cent; interest (on money)\textquotedblright{}.

\begin{quote}
\ru{Банк} \ru{даёт} \ru{пять} \ru{процентов} \ru{в} \ru{год}.
\end{quote}

\begin{quote}
\emph{Bank dayót pyat protséntov v god.}
\end{quote}

\begin{quote}
The bank pays five per cent a year.
\end{quote}

Masculine. \emph{\ru{Сто} \ru{процентов}} — a hundred per cent. \emph{\ru{Проценты}} is also interest on money.

\begin{cousinweb}
\textbf{\ru{процент}} is usually traced, through German \emph{Prozent}, to Latin \emph{prō centum} \textquotedblleft{}for a hundred\textquotedblright{} (Lewis \& Short).
\end{cousinweb}

\begin{grammarlens}[title={What you've built}]
One more word for telling, explaining and describing at B1.
\end{grammarlens}

\subsection*{Your turn}
\begin{itemize}
\raggedright
  \item \emph{Say it:} \emph{protsént}
  \item \emph{Say it:} \emph{protsént}, once more
  \item \emph{Say it:} a sentence of your own with \emph{protsént}
  \item \emph{Recall:} say the Russian for to lend, then the Russian for to give back, then say \emph{protsént} again
\end{itemize}

\begin{tcolorbox}[breakable,colback=teal!4,colframe=teal!35!black,title={Before you move on}]
What does \ru{процент} mean? (\textquotedblleft{}A per cent; interest (on money)\textquotedblright{}.) Say it once more.
\end{tcolorbox}
